% Lösungen Einheit 2 von 3 – Achsensymmetrie und Spiegeln (zusammenbau v0.1)
\einheitenkopf{Einheit 2 von 3 · Achsensymmetrie und Spiegeln}

\erg{16}{a) $2$ Achsen: die beiden Mittellinien}

16a)

\rechteck[achsen=beide]{5}{2}

\erg{17}{a) ja \quad b) die senkrechte Diagonale durch die obere und die untere Spitze \quad c) $1$ (Drachenviereck) \quad d) $2$ (die beiden Diagonalen) \quad e) Parallelogramm \quad f) $4$, $2$, $2$, $1$ (Quadrat, Rechteck, Raute, Drachenviereck) \quad g) $3$ \quad h) $2$ (senkrecht und waagerecht) \quad i) $AF = 6 : 2 = 3$\,cm \quad j) 1 \quad k) 1}

17b)

\drachen[achsen=senkrecht]{4}{3}{0.3}

\erg{18}{a) ja \quad b) $P'(6|3)$ \quad c) Bilddreieck $A'(7|1)$, $B'(5|1)$, $C'(6|4)$ \quad d) Bilddreieck $A'(1|7)$, $B'(4|7)$, $C'(4|5)$ \quad e) $A'(-2|1)$, $B'(-5|1)$, $C'(-4|4)$ \quad f) ergänzte Ecken $B'(6|1)$ und $C'(6|4)$; es entsteht ein Fünfeck in Hausform \quad g) Bilddreieck $A'(3|1)$, $B'(5|2)$, $C'(5|1)$ \quad h) ein gleichschenkliges Dreieck \quad i) Drachenviereck (a ist eine Diagonale, die andere Diagonale steht senkrecht darauf und wird halbiert; keine Raute, weil das Dreieck nicht gleichschenklig ist)}
\erg{19}{a) nein}
\erg{20}{a) Quadrat, Rechteck, Raute}
\erg{21}{a) Nur $1$ Achse: die waagerechte passt nicht, weil die beiden parallelen Seiten verschieden lang sind. \quad b) Faltet man es an einer Diagonale oder einer Mittellinie, passen die Hälften nicht aufeinander – die Ecken liegen versetzt. \quad c) z. B. ein Drachenviereck (Achse: eine Diagonale) oder ein gleichschenkliges Trapez \quad d) Rechteck}
